\documentclass[11pt]{beamer}

\usetheme{Madrid}
\usecolortheme{orchid}
\setbeamertemplate{navigation symbols}{}
\setbeamercolor{structure}{fg=orange!80!black}

\usepackage[utf8]{inputenc}
\usepackage[T1]{fontenc}
\usepackage[french]{babel}
\usepackage{amsmath}
\usepackage{booktabs}
\usepackage{array}
\usepackage{tikz}
\usetikzlibrary{arrows.meta, positioning, shapes.geometric}

\title{Sécurité de l'IA Générative}
\subtitle{Sprint 2 --- Étude architecturale et déduction des attaques\\[2pt]
\textcolor{orange!80!black}{\large Périmètre Orange : modèles discriminatifs / séquentiels}}
\author{Fatima-Zahra Rhoulam}
\institute{Télécom SudParis}
\date{}

\begin{document}

% ---------------------------------------------------------
\begin{frame}
  \titlepage
\end{frame}

% ---------------------------------------------------------
\begin{frame}{Plan de ma partie}
  \begin{enumerate}
    \item CNN de classification (CIFAR-10)
    \begin{itemize}
      \item Mécanisme exact (forward pass)
      \item Signal exploitable
      \item Attaque déduite : Membership Inference
    \end{itemize}
    \item Séries temporelles (LSTM)
    \begin{itemize}
      \item Mécanisme exact (portes, mémoire)
      \item Signal exploitable
      \item Attaque déduite : canaris \& exposure
    \end{itemize}
    \item Matrice Modèles $\leftrightarrow$ Attaques
  \end{enumerate}
\end{frame}

% ===========================================================
\section{CNN de classification (CIFAR-10)}

\begin{frame}{CNN --- le forward pass}
  \begin{center}
  \begin{tikzpicture}[
    box/.style={draw, rounded corners, fill=orange!10, minimum height=1cm, align=center, font=\footnotesize},
    arr/.style={-{Stealth[length=2mm]}, thick}
  ]
    \node[box] (img) {Image\\32$\times$32$\times$3};
    \node[box, right=0.5cm of img] (conv) {Convolution\\+ ReLU};
    \node[box, right=0.5cm of conv] (pool) {Pooling};
    \node[box, right=0.5cm of pool] (dense) {Flatten +\\Dense};
    \node[box, right=0.5cm of dense] (logits) {Logits\\(10 valeurs)};
    \node[box, right=0.5cm of logits] (soft) {Softmax};

    \draw[arr] (img) -- (conv);
    \draw[arr] (conv) -- (pool);
    \draw[arr] (pool) -- (dense);
    \draw[arr] (dense) -- (logits);
    \draw[arr] (logits) -- (soft);
  \end{tikzpicture}
  \end{center}
  \vspace{4mm}
  \begin{itemize}
    \item Convolution : filtres qui détectent des motifs locaux (contours, textures)
    \item Pooling : réduit la résolution, robustesse aux translations
    \item Softmax$(z_i) = \dfrac{\exp(z_i)}{\sum_j \exp(z_j)}$ $\rightarrow$ probabilité par classe
  \end{itemize}
\end{frame}

\begin{frame}{CNN --- la Loss comme origine du signal}
  \[
    \text{Loss} = -\log\big(p_{\text{classe correcte}}\big)
  \]
  \begin{itemize}
    \item La descente de gradient ajuste les poids pour minimiser cette Loss
          \textbf{sur les exemples d'entraînement}
    \item Sur une image d'entraînement : $p_{\max}$ très proche de 1 $\Rightarrow$ Loss $\approx 0$
    \item Sur une image de test (jamais vue) : $p_{\max}$ plus modérée $\Rightarrow$ Loss plus élevée
  \end{itemize}
  \begin{alertblock}{Conséquence}
    Cet écart de confiance train/test est la définition opérationnelle
    du surapprentissage --- et le signal exploité par la MIA.
  \end{alertblock}
\end{frame}

\begin{frame}{CNN --- attaque déduite : Membership Inference}
  \begin{center}
  \begin{tikzpicture}[
    box/.style={draw, rounded corners, fill=blue!8, minimum width=3.6cm, minimum height=1.1cm, align=center, font=\footnotesize},
    arr/.style={-{Stealth[length=2mm]}, thick}
  ]
    \node[box] (n1) {\textbf{1. Donnée cible} $x$\\IN ou OUT ?};
    \node[box, right=0.6cm of n1] (n2) {\textbf{2. Requête au CNN}\\récupère $p$ ou Loss};
    \node[box, below=0.6cm of n2] (n3) {\textbf{3. Signal observé}\\$p_{\max}$ ou Loss$=-\log p$};
    \node[box, left=0.6cm of n3] (n4) {\textbf{4. Seuil calibré}\\(shadow models / Yeom)};
    \node[box, below=0.6cm of n4] (n5) {\textbf{5. Verdict}\\IN si signal fort};

    \draw[arr] (n1) -- (n2);
    \draw[arr] (n2) -- (n3);
    \draw[arr] (n3) -- (n4);
    \draw[arr] (n4) -- (n5);
  \end{tikzpicture}
  \end{center}
  \begin{itemize}
    \item Shokri et al. (2017) : shadow models + classifieur d'attaque
    \item Yeom et al. (2018) : seuil direct sur la Loss, taux de succès
          borné par l'écart train/test accuracy
  \end{itemize}
\end{frame}

\begin{frame}{CNN --- exemple chiffré}
  \begin{table}
  \centering
  \begin{tabular}{lcc}
    \toprule
     & \textbf{Image d'entraînement} & \textbf{Image de test} \\
    \midrule
    $p_{\max}$ (classe ``chat'') & 0{,}994 & 0{,}71 \\
    Loss $=-\log(p_{\max})$ & 0{,}006 & 0{,}342 \\
    \bottomrule
  \end{tabular}
  \end{table}
  \vspace{3mm}
  Avec un seuil Loss $< 0{,}05 \Rightarrow$ IN, les deux exemples
  sont classés correctement par l'attaquant.
\end{frame}

% ===========================================================
\section{Séries temporelles (LSTM)}

\begin{frame}{LSTM --- le forward pass}
  \begin{center}
  \begin{tikzpicture}[
    box/.style={draw, rounded corners, fill=orange!10, minimum width=2.6cm, minimum height=1cm, align=center, font=\footnotesize},
    arr/.style={-{Stealth[length=2mm]}, thick}
  ]
    \node[box] (f) {Porte d'oubli\\$f_t$};
    \node[box, right=0.5cm of f] (i) {Porte d'entrée\\$i_t$};
    \node[box, right=0.5cm of i] (c) {Mémoire\\$C_t$};
    \node[box, right=0.5cm of c] (o) {Porte de sortie\\$o_t$};
    \node[box, right=0.5cm of o] (pred) {Prédiction\\$\hat{x}_{t+1}$};

    \draw[arr] (f) -- (i);
    \draw[arr] (i) -- (c);
    \draw[arr] (c) -- (o);
    \draw[arr] (o) -- (pred);
  \end{tikzpicture}
  \end{center}
  \vspace{3mm}
  \[
    C_t = f_t \odot C_{t-1} + i_t \odot \tilde{C}_t
  \]
  \begin{itemize}
    \item $f_t$ : proportion de mémoire conservée
    \item $i_t$ : proportion de nouvelle information ajoutée
    \item $o_t$ : proportion de mémoire exposée dans $h_t \to \hat{x}_{t+1}$
  \end{itemize}
\end{frame}

\begin{frame}{LSTM --- l'erreur de prédiction comme signal}
  \[
    \text{Loss}_t = \big(\hat{x}_{t+1} - x_{t+1}\big)^2 \quad \text{(MSE)}
  \]
  \begin{itemize}
    \item Rétropropagation à travers le temps (BPTT)
    \item Séquence d'entraînement : erreur résiduelle anormalement basse
    \item Séquence de test : erreur plus élevée en moyenne
  \end{itemize}
  \begin{block}{Analogie}
    Équivalent, pour une série numérique continue, de la
    \textbf{perplexité} pour un LLM : une mesure de ``surprise''
    du modèle, ici en MSE plutôt qu'en probabilité de token.
  \end{block}
\end{frame}

\begin{frame}{LSTM --- attaque déduite : canaris \& exposure}
  \begin{center}
  \begin{tikzpicture}[
    box/.style={draw, rounded corners, fill=blue!8, minimum width=3.8cm, minimum height=1.1cm, align=center, font=\footnotesize},
    arr/.style={-{Stealth[length=2mm]}, thick}
  ]
    \node[box] (n1) {\textbf{1. Injection}\\canari rare, $N$ exemplaires};
    \node[box, right=0.6cm of n1] (n2) {\textbf{2. Entraînement}\\checkpoints 20/100/200 ép.};
    \node[box, below=0.6cm of n2] (n3) {\textbf{3. Requête}\\erreur sur le vrai canari};
    \node[box, left=0.6cm of n3] (n4) {\textbf{4. Comparaison}\\à toutes les variantes};
    \node[box, below=0.6cm of n4] (n5) {\textbf{5. Rang = exposure}\\Carlini et al. (2019)};

    \draw[arr] (n1) -- (n2);
    \draw[arr] (n2) -- (n3);
    \draw[arr] (n3) -- (n4);
    \draw[arr] (n4) -- (n5);
  \end{tikzpicture}
  \end{center}
  \begin{itemize}
    \item Méthode fondatrice : \emph{The Secret Sharer}, Carlini et al., USENIX Security 2019
    \item Testée à l'origine sur un LSTM de génération de texte
  \end{itemize}
\end{frame}

\begin{frame}{LSTM --- exemple chiffré}
  \begin{table}
  \centering
  \begin{tabular}{lc}
    \toprule
    \textbf{Mesure} & \textbf{Valeur} \\
    \midrule
    MSE sur le canari exact & 0{,}0004 \\
    MSE moyenne sur séquences de test & 0{,}183 \\
    Rang du canari parmi 10 variantes & 1$^{\text{er}}$/10 \\
    \bottomrule
  \end{tabular}
  \end{table}
  \vspace{3mm}
  En comparant les checkpoints (20 $\to$ 100 $\to$ 200 époques), le
  rang se dégrade progressivement : on visualise \textbf{quand}
  la mémorisation apparaît pendant l'entraînement.
\end{frame}

% ===========================================================
\section{Synthèse}

\begin{frame}{Matrice Modèles $\leftrightarrow$ Attaques --- contribution Orange}
  \begin{table}
  \centering
  \scriptsize
  \begin{tabular}{p{2.4cm}p{2.2cm}p{2.6cm}p{2.2cm}}
    \toprule
    \textbf{Modèle} & \textbf{Sortie} & \textbf{Signal exploité} & \textbf{Attaque} \\
    \midrule
    CNN (CIFAR-10) & Probabilité/classe & Écart Loss train/test & MIA (Shokri 2017, Yeom 2018) \\
    LSTM (TSGM) & Valeur continue & MSE anormalement basse + rang d'exposure & MIA + canaris (Carlini 2019) \\
    \bottomrule
  \end{tabular}
  \end{table}
\end{frame}

\begin{frame}{Argument transversal}
  \begin{alertblock}{Point clé}
    Dans les deux cas (CNN et LSTM), le modèle ne renvoie
    \textbf{jamais la donnée elle-même} --- contrairement au
    périmètre Bleu (génératif), où la sortie peut \emph{être}
    la donnée mémorisée.
  \end{alertblock}
  \begin{itemize}
    \item Côté Orange : uniquement des attaques d'\textbf{inférence}
          (verdict IN/OUT, score d'exposure gradué)
    \item Jamais d'extraction verbatim, contrairement à un LLM ou
          un modèle de diffusion
    \item Démonstration pratique de la distinction prédictif / génératif
  \end{itemize}
\end{frame}

\begin{frame}{Sources}
  \footnotesize
  \begin{itemize}
    \item Shokri, Stronati, Song \& Shmatikov (2017). \emph{Membership
          Inference Attacks Against Machine Learning Models}. IEEE S\&P.
    \item Yeom, Giacomelli, Fredrikson \& Jha (2018). \emph{Privacy Risk
          in Machine Learning: Analyzing the Connection to Overfitting}.
          IEEE CSF. \url{https://arxiv.org/abs/1709.01604}
    \item Carlini, Liu, Erlingsson, Kos \& Song (2019). \emph{The Secret
          Sharer: Evaluating and Testing Unintended Memorization in
          Neural Networks}. USENIX Security.
          \url{https://arxiv.org/abs/1802.08232}
    \item Fredrikson, Jha \& Ristenpart (2015). \emph{Model Inversion
          Attacks that Exploit Confidence Information and Basic
          Countermeasures}. ACM CCS.
  \end{itemize}
\end{frame}

\end{document}
